\documentclass[12pt, titlepage]{article}

\usepackage{amsmath, mathtools}

\usepackage[round]{natbib}
\usepackage{amsfonts}
\usepackage{amssymb}
\usepackage{graphicx}
\usepackage{colortbl}
\usepackage{xr}
\usepackage{hyperref}
\usepackage{longtable}
\usepackage{xfrac}
\usepackage{tabularx}
\usepackage{float}
\usepackage{siunitx}
\usepackage{booktabs}
\usepackage{multirow}
\usepackage[section]{placeins}
\usepackage{caption}
\usepackage{fullpage}

\hypersetup{
bookmarks=true,     % show bookmarks bar?
colorlinks=true,       % false: boxed links; true: colored links
linkcolor=red,          % color of internal links (change box color with linkbordercolor)
citecolor=blue,      % color of links to bibliography
filecolor=magenta,  % color of file links
urlcolor=cyan          % color of external links
}

\usepackage{array}


\input{../../Comments.text}
\input{../../Common.text}

\begin{document}

\title{Module Interface Specification for \progname{}}

\author{\authname}

\date{\today}

\maketitle

\pagenumbering{roman}

\section{Revision History}

\begin{tabularx}{\textwidth}{p{3cm}p{2cm}X}
\toprule {\bf Date} & {\bf Version} & {\bf Notes}\\
\midrule
Nov 7, 2026 & 1.0 & Davis Lenover: Revision $-1$. 4 Notation, and the MIS of
Database, Clock, Market Providers, Control, Constants, Market, Trading, Plugins
and JSON Format.\\
Nov 9, 2026 & 1.1 & Davis Lenover: the MIS of Users and Wallet, with the reason
for each decision. 1 to 5: Revision History, Symbols, Introduction, Notation and
Module Decomposition.\\
\bottomrule
\end{tabularx}

~\newpage

\section{Symbols, Abbreviations and Acronyms}

This document uses the terms and acronyms of the
\href{https://github.com/davislenover/BurnNEarn/blob/main/pdfs/SRS-IEEE/SRS.pdf}{SRS}
(SRS Section~1.4) and its symbolic constants, such as MAX\_DAILY\_CURRENCY
(SRS Appendix~C). Constants (M7) gives the code their values, which must match
the SRS. The abbreviations this document uses are:

\begin{center}
\renewcommand{\arraystretch}{1.2}
\begin{tabular}{l l}
\toprule
\textbf{Symbol} & \textbf{Description}\\
\midrule
API & Application Programming Interface\\
FR & Functional Requirement (SRS), by its printed number\\
H & Hazard (HA), such as H51\\
HA & Hazard Analysis\\
HTTP & Hypertext Transfer Protocol\\
IANA & Internet Assigned Numbers Authority, which names time zones\\
ID & Identifier\\
JSON & JavaScript Object Notation\\
M & Module (MG), such as M8\\
MG & Module Guide\\
MIS & Module Interface Specification\\
NaN & Not a Number, a Decimal value that is not a number\\
NFR & Nonfunctional Requirement (SRS)\\
SEC & Security Requirement (SRS and HA)\\
SRS & Software Requirements Specification\\
UC & Use Case (SRS)\\
URL & Uniform Resource Locator\\
USD & United States dollar\\
UTC & Coordinated Universal Time\\
UUID & Universally Unique Identifier\\
\bottomrule
\end{tabular}
\end{center}

\newpage

\tableofcontents

\newpage

\pagenumbering{arabic}

\section{Introduction}

This document is the Module Interface Specification (MIS) for \progname{}, a
phone app where physical activity earns virtual currency that the user invests
in a simulated stock market. \progname{} is two programs: a server, written in
Python with FastAPI and backed by PostgreSQL, and a phone app, written in React
Native. The Development Plan (Section~10) gives the reason for each of these
choices.

The MIS gives each module's interface: the types, constants and access programs
it exports, and what each access program does. The Module Guide (MG) says what
each module hides and why the system is split this way. The Software
Requirements Specification (SRS) and the Hazard Analysis (HA) say what the
system must do. Next to each decision, a module's section gives its reason: the
requirement or hazard it serves, or the design pattern it uses. A later change
can then be checked against what the decision was for.

This revision specifies the modules of the proof of concept. Modules marked
\textsuperscript{\dag} in Table~\ref{TblMH} are built after it and are
specified at Revision~0, so that every section describes an interface the code
really has.

The full documentation and implementation are in the project repository:
\url{https://github.com/davislenover/BurnNEarn}.

\newpage

\section{Notation}

The structure of the MIS for modules comes from \citet{HoffmanAndStrooper1995},
with the addition that template modules have been adapted from
\cite{GhezziEtAl2003}.  The mathematical notation comes from Chapter 3 of
\citet{HoffmanAndStrooper1995}.  For instance, the symbol := is used for a
multiple assignment statement and conditional rules follow the form $(c_1
\Rightarrow r_1 | c_2 \Rightarrow r_2 | ... | c_n \Rightarrow r_n )$. The
table below gives the primitive data types.

\begin{center}
\renewcommand{\arraystretch}{1.2}
\noindent 
\begin{tabular}{l l p{9.6cm}} 
\toprule 
\textbf{Data Type} & \textbf{Notation} & \textbf{Description}\\ 
\midrule
integer & $\mathbb{Z}$ & a number without a fractional component in (-$\infty$, $\infty$) \\
natural number & $\mathbb{N}$ & a number without a fractional component in [1, $\infty$) \\
real & $\mathbb{R}$ & any number in (-$\infty$, $\infty$)\\
boolean & $\mathbb{B}$ & a value that is either true or false\\
string & String & a sequence of characters\\
decimal & Decimal & an exact decimal number, such as a price, so money is never rounded by accident (NFR1)\\
date & Date & a calendar day, without a time of day\\
date and time & DateTime & a moment in time with its time zone, so the server, the phone and the market agree on it (FR33, SEC-24)\\
duration & Duration & a length of time, such as 15 minutes\\
identifier & ID & a UUID, which the phone can make without the database, such as a trade's request ID (FR38)\\
\bottomrule
\end{tabular} 
\end{center}

\noindent \progname{} also uses derived types. A sequence of values of one type
T is ``seq of T'', a set is ``set of T'', and a tuple, whose values may have
different types, is ``tuple of (name: T, ...)''. A string is a sequence of
characters, and a value that may be missing is ``T or None''. A function from
T1 to T2 is T1 $\rightarrow$ T2, and a local function is given by its type
signature and then its specification. ``new X'' or X($v$, ...) makes a value
of type X, $t$.f is the value of field f in tuple $t$, and $p$.f($x$) calls
access program f of $p$, which the In column leaves out. A type or a constant
is defined with $=$, a type written \{a, b\} has just the values a and b, and
$:=$ gives a value in the semantics. There,
transition says how the state changes, output gives $out$, the value returned,
and exception gives $exc$, the exception raised. When no case of $exc$ holds,
nothing is raised.

Types such as Session, Plugin and DomainError are exported by a module and
defined in its section. ValueError is Python's, IntegrityError and
NoResultFound are SQLAlchemy's, and ValidationError is Pydantic's. In
an access program table, In gives each input's type in order, separated by
commas, and - means none. A route is written as its HTTP method and path, such
as GET /v1/me/balance. Its In column lists what the phone
sends, and its Exceptions column the refusals it raises, besides the 422 that
JSON Format (M23) describes for a request that doesn't fit the route's types
(SEC-25).

\newpage

\section{Module Decomposition}

The following table is taken directly from the MG. Each section of this
document keeps its module's MG number, such as Users (M8), so a module can be
traced between the two documents. Modules marked (app) are in the phone app,
Constants is in both programs, and the rest are in the server.

\begin{table}[h!]
\centering
\begin{tabular}{p{0.3\textwidth} p{0.6\textwidth}}
\toprule
\textbf{Level 1} & \textbf{Level 2}\\
\midrule

\multirow{6}{0.3\textwidth}{Hardware-Hiding Module} & M1 Database\\
& M2 Clock\\
& M3 Market Providers\\
& M25 Email Service\textsuperscript{\dag}\\
& M4 Health Service (app)\\
& M5 Device Service (app)\\
\midrule

\multirow{16}{0.3\textwidth}{Behaviour-Hiding Module} & M6 Control\\
& M7 Constants (server and app)\\
& M8 Users\\
& M9 Wallet\\
& M10 Market\\
& M11 Activity\\
& M12 Trading\\
& M13 Goals\textsuperscript{\dag}\\
& M14 Learning\textsuperscript{\dag}\\
& M15 Social\textsuperscript{\dag}\\
& M16 Scenarios\textsuperscript{\dag}\\
& M17 Privacy\textsuperscript{\dag}\\
& M18 Navigation (app)\\
& M19 Screens (app)\\
& M20 Sync Service (app)\\
& M21 Reminder Service (app)\textsuperscript{\dag}\\
\midrule

\multirow{3}{0.3\textwidth}{Software Decision Module} & M22 Plugins\\
& M23 JSON Format\\
& M24 API Service (app)\\
\bottomrule

\end{tabular}
\caption{Module Hierarchy}
\label{TblMH}
\end{table}

\noindent Modules marked \textsuperscript{\dag} are built after the proof of
concept and are specified at Revision~0.

\newpage

\section{MIS of Database} \label{MisDB}

\subsection{Module}

Database (M1 in the MG)

\subsection{Uses}

None

\subsection{Syntax}

\subsubsection{Exported Types}

\noindent Session: one unit of work on the database, saved by commit and
dropped by rollback or close.

\medskip

\noindent Base: the type every module defines its own tables on. Every
timestamp keeps its time zone.

\medskip
\noindent A table holds one kind of thing, such as assets. It has a row for
each asset and a column for each fact about it, such as its symbol.

\medskip
\noindent A table can have constraints. A constraint is a rule the database
itself checks every time a row is added, changed or deleted, and it refuses a
change that breaks the rule. There are five kinds, shown here with examples
from Market (M10):
\begin{itemize}
\item Not null: every row must have a value in this column. Every column in
  Market is not null.
\item Primary key: the column that tells each row apart, so no two rows have
  the same value in it. A table has one. Each asset has its own asset\_id.
\item Foreign key: a column that must hold the primary key of a row in another
  table. Each price's asset\_id must be the asset\_id of an asset, so a price
  can never belong to an asset that does not exist. For the same reason, an
  asset that still has prices cannot be deleted.
\item Unique: no two rows may have the same value in the columns it names. A
  table can have several. No two assets share a symbol. With more than one
  column, only the combination must differ: an asset has many prices, but at
  most one for each quoted\_at (the time the price was set).
\item Check: a condition every row must meet, such as a price being above 0.
\end{itemize}
Last, a table can have an index. Like the index of a book, it lets the
database find rows by a column without reading every row. An index is not a
rule. The table of calls to the market data provider, provider\_requests, has
one on requested\_at, so counting the calls in the last minute stays fast. A
primary key or unique comes with its own index, so it needs no second one.

\medskip
\noindent Base names every primary key, foreign key, unique, check and index
that has no name of its own, from its table and columns, by the rules below,
instead of letting the database choose. A migration (a script that changes the
tables in the database) can then change or drop one by a name known when the
migration is written:

\begin{center}
\begin{tabular}{p{2.4cm} p{5.6cm} p{6.6cm}}
\toprule
\textbf{Kind} & \textbf{Name} & \textbf{Example} \\
\midrule
Primary key & pk\_\textit{table} & pk\_assets \\
Foreign key & fk\_\textit{table}\_\textit{first column}\_\textit{other table} & fk\_prices\_asset\_id\_assets \\
Unique & uq\_\textit{table}\_\textit{each column} & uq\_prices\_asset\_id\_quoted\_at \\
Check & ck\_\textit{table}\_\textit{check name} & ck\_prices\_price\_positive \\
Index & ix\_\textit{table}\_\textit{first column} & ix\_provider\_requests\_requested\_at \\
\bottomrule
\end{tabular}
\end{center}

\noindent A check is a condition, not a list of columns, so whoever writes it
gives it a name, such as price\_positive, and Base adds the rest. A check
without a name is an error.

\subsubsection{Exported Constants}

None

\subsubsection{Exported Access Programs}

\begin{center}
\begin{tabular}{p{3cm} p{2cm} p{3cm} p{3cm}}
\toprule
\textbf{Name} & \textbf{In} & \textbf{Out} & \textbf{Exceptions} \\
\midrule
SessionLocal & - & Session & - \\
get\_db & - & Session & - \\
\bottomrule
\end{tabular}
\end{center}

\subsection{Semantics}

\subsubsection{State Variables}

engine: the connection to the database at DATABASE\_URL, made once when the
module is first loaded.

\subsubsection{Environment Variables}

\noindent database: the PostgreSQL server that holds every module's tables.

\medskip

\noindent DATABASE\_URL: String, the database to connect to, read from the
environment or .env.

\subsubsection{Assumptions}

\begin{itemize}
\item The migrations are applied before the server starts.
\item Each module owns its tables. It never reads or writes another module's
  tables, though it may have a foreign key into them.
\item Only routes commit, at most once per request. A route is the function in
  a module's routes.py that answers one HTTP request, such as GET /v1/assets in
  Market (M10). To commit is to save a session's changes to the database for
  good.
\end{itemize}

\noindent The commit rule works like this. The code a route calls, in its own
module or another, changes the session but never commits. The route commits
once, after every step has succeeded. If a step refuses the request or fails,
the route does not commit, so nothing is saved. So everything one request
changes, such as a buy's payment and its new shares, is saved together or not
at all. The price poller in Market (M10) is the one exception: it answers no
request, so it commits once per poll.

\subsubsection{Access Routine Semantics}

\noindent SessionLocal():
\begin{itemize}
\item output: $out :=$ a new Session on engine. The caller closes it, and
  commits only as the Assumptions allow.
\item exception: None
\end{itemize}

\noindent get\_db():
\begin{itemize}
\item transition: when the request ends, the session is closed and anything
  the route did not commit is dropped.
\item output: $out :=$ SessionLocal(), a session for one request alone. Routes
  ask for it as a dependency, so a test can replace get\_db with a session whose
  changes are undone after the test.
\item exception: None. The session raises OperationalError, a kind of
  SQLAlchemyError, at its first use if the database cannot be reached.
\end{itemize}

\subsubsection{Local Functions}

None

\newpage

\section{MIS of Clock} \label{MisClock}

\subsection{Module}

Clock (M2 in the MG)

\subsection{Uses}

None

\subsection{Syntax}

\subsubsection{Exported Types}

Clock = () $\rightarrow$ DateTime, a function that returns the current time.

\subsubsection{Exported Constants}

None

\subsubsection{Exported Access Programs}

\begin{center}
\begin{tabular}{p{3cm} p{2cm} p{3cm} p{3cm}}
\toprule
\textbf{Name} & \textbf{In} & \textbf{Out} & \textbf{Exceptions} \\
\midrule
system\_clock & - & DateTime & - \\
get\_clock & - & Clock & - \\
\bottomrule
\end{tabular}
\end{center}

\subsection{Semantics}

\subsubsection{State Variables}

None

\subsubsection{Environment Variables}

osTime: DateTime, the current time from the server's operating system clock.

\subsubsection{Assumptions}

The server's operating system clock has these properties:
\begin{itemize}
\item It is synchronized to Coordinated Universal Time (UTC) by the Network
  Time Protocol, and stays within 1 second of UTC. That is close enough, since
  market hours, the price staleness limit and the server's date (SEC-24 in the
  SRS) all work in whole minutes or more.
\item It is not monotonic. A sync can move it back slightly, so a later call
  can return an earlier time. That does no harm. Its times are stored and
  compared with other stored times and with market hours, and nothing uses it
  to time a wait.
\end{itemize}

\subsubsection{Access Routine Semantics}

\noindent system\_clock():
\begin{itemize}
\item output: $out := $ osTime in UTC, to the microsecond
\item exception: None
\end{itemize}

\noindent get\_clock():
\begin{itemize}
\item output: $out := $ system\_clock. The output is the function itself, not
  a time. Routes call it to read the current time, so a test can replace
  get\_clock with a Clock that always returns a fixed time.
\item exception: None
\end{itemize}

\subsubsection{Local Functions}

None

\newpage

% The sections follow the MG's module order, and each starts on a new page.
\section{MIS of Market Providers} \label{MisProviders}

\subsection{Module}

Market Providers (M3 in the MG)

\subsection{Uses}

Clock (M2)

\subsection{Syntax}

\subsubsection{Exported Types}

\noindent Quote = tuple of (symbol: String, price: Decimal, quoted\_at:
DateTime), a stock's latest price and when the provider set it (not when the
server asked for it).

\medskip
\noindent MarketDataProvider: the abstract type of every provider. It has one
access program, get\_quote. The providers are currently FinnhubProvider and
FakeProvider.

\medskip
\noindent ProviderSettings = tuple of (market\_provider: String,
finnhub\_api\_key: String or None), the name of the provider to use,
currently ``finnhub'' or ``fake'', and the key for Finnhub.

\medskip
\noindent MarketDataError, the error for no usable price, with its kinds
RateLimitedError (the request limit was hit) and UnknownSymbolError (no price
for the symbol).

\subsubsection{Exported Constants}

FINNHUB\_CALLS\_PER\_MINUTE: $\mathbb{N}$ = 60, Finnhub's free plan limit.

\subsubsection{Exported Access Programs}

\begin{center}
\begin{tabular}{>{\raggedright\arraybackslash}p{3.95cm} p{3.05cm} p{3.75cm} p{4.05cm}}
\toprule
\textbf{Name} & \textbf{In} & \textbf{Out} & \textbf{Exceptions} \\
\midrule
new FinnhubProvider & String & FinnhubProvider & - \\
new FakeProvider & Clock & FakeProvider & - \\
get\_quote & String & Quote & MarketDataError,\newline RateLimitedError,\newline
UnknownSymbolError \\
build\_provider & ProviderSettings & MarketDataProvider & ValueError \\
get\_provider\_settings & - & ProviderSettings & ValidationError \\
\bottomrule
\end{tabular}
\end{center}

\subsection{Semantics}

\subsubsection{State Variables}

\noindent key: String, a FinnhubProvider's API key, sent only to Finnhub in a
request header. It is never put in a URL or an error message, and never sent
to the phone.\\
clock: Clock and prices: String $\rightarrow$ Decimal, a FakeProvider's clock
and last prices.

\subsubsection{Environment Variables}

\noindent finnhubApi: the Finnhub quote API, whose reply has a status, a price
$c$ and its time $t$.

\medskip
\noindent MARKET\_PROVIDER, FINNHUB\_API\_KEY: String, the two fields of
ProviderSettings, read from the environment or the .env file.

\subsubsection{Assumptions}

\begin{itemize}
\item Every provider follows MarketDataProvider: its get\_quote returns a Quote
  or raises a MarketDataError, never another error.
\item build\_provider is only given settings from get\_provider\_settings, so
  they hold all that the named provider needs.
\end{itemize}

\subsubsection{Access Routine Semantics}

\noindent new FinnhubProvider($k$), new FakeProvider($f$):
\begin{itemize}
\item output: $out :=$ FinnhubProvider(key $= k$) or FakeProvider(clock $= f$,
  no prices)
\item exception: None
\end{itemize}

\noindent $p$.get\_quote($s$), for a FinnhubProvider $p$:
\begin{itemize}
\item output: $out :=$ (symbol: $s$, price: $c$, quoted\_at: $t$ seconds after
  1970-01-01 00:00 UTC), from one request to finnhubApi holding only $s$ and
  key (10 second timeout)
\item exception: $exc :=$ (the request fails or its reply cannot be read
  $\Rightarrow$ \mbox{MarketDataError} $|$ status $= 429 \Rightarrow$ RateLimitedError
  $|$ status $\neq 200$, or no finite $c$ or valid $t$ $\Rightarrow$
  \mbox{MarketDataError} $|$ $c \leq 0 \lor t = 0 \Rightarrow$ UnknownSymbolError), in
  that order
\end{itemize}

\noindent $p$.get\_quote($s$), for a FakeProvider $p$:
\begin{itemize}
\item transition: prices($s$) $:=$ (no price for $s$ yet $\Rightarrow$
  start($s$) $|$ otherwise $\Rightarrow$ step(prices($s$)))
\item output: $out :=$ (symbol: $s$, price: prices($s$), quoted\_at: clock())
\item exception: None, even for a symbol no real stock has
\end{itemize}

\noindent build\_provider($v$):
\begin{itemize}
\item output: $out :=$ the provider that $v$ names. That is currently new
  FinnhubProvider with $v$'s key for ``finnhub'', or new
  FakeProvider(system\_clock) for ``fake''.
\item exception: $exc :=$ ($v$.market\_provider names no provider
  $\Rightarrow$ ValueError)
\end{itemize}

\noindent get\_provider\_settings():
\begin{itemize}
\item output: $out :=$ the settings in MARKET\_PROVIDER and FINNHUB\_API\_KEY,
  read once, on the first call that succeeds. Unset, they give ``fake'' and None.
\item exception: $exc :=$ (MARKET\_PROVIDER names no provider, or names one
  that needs a key, currently only Finnhub, and the key is unset or empty
  $\Rightarrow$ ValidationError)
\end{itemize}

\subsubsection{Local Functions}

\noindent start: String $\rightarrow$ Decimal, start($s$) $:=$ a price from 20
to 500 that depends only on $s$\\
step: Decimal $\rightarrow$ Decimal, step($x$) $:=$ $x$ moved up or down at
random by at most 0.5\%, to four decimals

\newpage

% M4 Health Service and M5 Device Service (app): #250
\section{MIS of Control} \label{MisControl}

\subsection{Module}

Control (M6 in the MG)

\subsection{Uses}

Database (M1), Users (M8), Wallet (M9), Market (M10), Activity (M11), Trading
(M12), Plugins (M22), JSON Format (M23)

\medskip
\noindent For now PLUGINS holds only Market (M10). The others are added as they
are built.

\subsection{Syntax}

\subsubsection{Exported Types}

None

\subsubsection{Exported Constants}

\noindent PLUGINS: seq of Plugin $=$ $\langle$PLUGIN of Market (M10)$\rangle$,
the PLUGIN of each module in the server, in the order their background work
starts. Listing a module here is also how database migrations find its tables.

\subsubsection{Exported Access Programs}

\begin{center}
\begin{tabular}{p{3cm} p{2.4cm} >{\raggedright\arraybackslash}p{6.8cm} p{2.3cm}}
\toprule
\textbf{Name} & \textbf{In} & \textbf{Out} & \textbf{Exceptions} \\
\midrule
create\_app & seq of Plugin & App & - \\
GET /v1/health & - & tuple of (status: String, database: String, modules:
  String $\rightarrow$ ModuleHealth) & - \\
\bottomrule
\end{tabular}
\end{center}

\subsection{Semantics}

\subsubsection{State Variables}

app: App, the running server, built by create\_app(PLUGINS) when
\texttt{uvicorn app.main:app} starts the server.

\subsubsection{Environment Variables}

None

\subsubsection{Assumptions}

No two plug-ins in PLUGINS share a name, and no two routes share a method and
path. Each plug-in's lifespan and health check follow the rules in Plugins
(M22).

\subsubsection{Access Routine Semantics}

\noindent create\_app(\textit{plugins}):
\begin{itemize}
\item output: $out :=$ a new App that serves GET /v1/health, which reports
  the health of each plug-in in \textit{plugins}, and the routes of
  each plug-in in \textit{plugins}, all under /v1, plus the API docs at /docs.
  It answers any DomainError a route raises with handle\_domain\_error, and
  runs lifespan when it starts and shuts down. \textit{plugins} is PLUGINS when
  none are given.
\item exception: None
\end{itemize}

\noindent GET /v1/health():
\begin{itemize}
\item output: $out :=$ (status: $s$, database: $d$, modules: $m$), where
  \begin{itemize}
  \item $d :=$ (a query to the database succeeds $\Rightarrow$ ``ok'' $|$ it
    fails $\Rightarrow$ ``unavailable'')
  \item for each plug-in $p$ of the App with a health check, $m(p$.name$) :=$
    ($p$.health() raises an error $\Rightarrow$ (\mbox{``degraded''}, ``health
    check failed''), and the error is logged $|$ otherwise $\Rightarrow$
    $p$.health()). A plug-in without a health check is left out of $m$.
  \item $s :=$ ($d =$ ``ok'' $\land$ every status in $m$ is ``ok''
    $\Rightarrow$ ``ok'' $|$ otherwise $\Rightarrow$ ``degraded'')
  \end{itemize}
  The reply has HTTP status 503 when $d$ is ``unavailable'', since nothing
  works without the database, and 200 otherwise, even when a module is
  degraded.
\item exception: None
\end{itemize}

\subsubsection{Local Functions}

\noindent lifespan: Lifespan, the background work of each App that create\_app
builds, made from its \textit{plugins}.

\medskip
\noindent lifespan($a$):
\begin{itemize}
\item transition: the lifespans of \textit{plugins} start in order when $a$
  starts and stop in reverse order when $a$ shuts down, and if one fails to
  start, the ones already started still stop.
\item exception: $exc :=$ (a lifespan raises an error $\mathit{err}$ while it
  starts $\Rightarrow$ $\mathit{err}$). $a$ does not start.
\end{itemize}

\medskip
\noindent handle\_domain\_error: DomainError $\rightarrow$ tuple of (status:
$\mathbb{N}$, body: Refusal), the reply sent when a route raises a DomainError.

\medskip
\noindent handle\_domain\_error($e$):
\begin{itemize}
\item output: $out :=$ (status: $e$.status, body: the Refusal of $e$ in JSON
  Format (M23))
\item exception: None
\end{itemize}

\newpage

\section{MIS of Constants} \label{MisConstants}

\subsection{Module}

Constants (M7 in the MG)

\subsection{Uses}

None

\subsection{Syntax}

\subsubsection{Exported Types}

{\raggedright
ActivityType = \{STEPS, WHEELCHAIR\_PUSHES, CYCLING, SWIMMING, OTHER\_WORKOUT\},
the supported activity types (SUPPORTED\_ACTIVITY\_TYPES in SRS Appendix C).
Each value is also a String, its name in lower case, such as ``steps''.\par}

\subsubsection{Exported Constants}

Names and values match SRS Appendix C. EARNING\_RATES($t$) is the amount of
activity type $t$ that earns one unit of virtual currency ($u_t$ in SRS
Section~6.7). Decimal amounts are in virtual currency. The phone app's copy,
once written, must hold the same values.

\begin{longtable}{p{6cm} p{3.4cm} >{\raggedright\arraybackslash}p{5.4cm}}
\toprule
\textbf{Name} & \textbf{Type} & \textbf{Value} \\
\midrule
\endhead
\multicolumn{3}{l}{\textit{Activity and rewards}} \\
EARNING\_RATES & ActivityType $\rightarrow \mathbb{N}$ & 100 steps for STEPS,
  100 pushes for WHEELCHAIR\_PUSHES, 1 minute for each other type \\
MAX\_DAILY\_CURRENCY & Decimal & 300 \\
MAX\_DAILY\_STEPS & $\mathbb{N}$ & 150,000 \\
MAX\_DAILY\_PUSHES & $\mathbb{N}$ & 100,000 \\
MAX\_WORKOUT\_DURATION & Duration & 12 hours \\
SYNC\_LOOKBACK\_PERIOD & Duration & 7 days \\
\midrule
\multicolumn{3}{l}{\textit{Simulated market}} \\
QUOTE\_CURRENCY & String & ``USD'' (US dollars) \\
QUOTE\_STALENESS\_LIMIT & Duration & 15 minutes \\
MIN\_TRADABLE\_STOCKS & $\mathbb{N}$ & 50 \\
SCENARIO\_START\_BALANCE & Decimal & 10,000 \\
VALUE\_TOLERANCE & Decimal & 0.01 \\
\midrule
\multicolumn{3}{l}{\textit{Accounts and data}} \\
MIN\_USER\_AGE & $\mathbb{N}$ & 18 years \\
MIN\_PASSWORD\_LENGTH & $\mathbb{N}$ & 8 characters \\
MAX\_FAILED\_SIGNINS & $\mathbb{N}$ & 5 \\
SIGNIN\_LOCKOUT\_PERIOD & Duration & 15 minutes \\
SESSION\_TIMEOUT & Duration & 14 days \\
DATA\_DELETION\_PERIOD & Duration & 30 days \\
\bottomrule
\end{longtable}

\subsubsection{Exported Access Programs}

None

\subsection{Semantics}

\subsubsection{State Variables}

None

\subsubsection{Environment Variables}

None

\subsubsection{Assumptions}

None

\subsubsection{Access Routine Semantics}

None

\subsubsection{Local Functions}

None

\newpage

% M8 Users and M9 Wallet: #248
\section{MIS of Users} \label{MisUsers}

\subsection{Module}

Users (M8 in the MG)

\subsection{Uses}

Database (M1), Plugins (M22), JSON Format (M23). The MG also lists Clock (M2),
Constants (M7) and Email Service (M25), for sign-up, sign-in, sessions and the
password reset link (FR6) at Revision~0.

\subsection{Syntax}

\subsubsection{Exported Types}

\noindent CurrentUser = tuple of (user\_id: ID), a read-only copy of the user a
request comes from. Routes get it from get\_current\_user, never from the
request, so the phone can't act as another user (SEC-15).

\medskip
\noindent NotAuthenticatedError, a DomainError with code
``not\_authenticated'', status 401 and the message ``Sign in to continue.''
It refuses a request with no signed-in user (SEC-15). Status 401 is HTTP's
status for that, so the phone can tell it apart from other refusals.

\subsubsection{Exported Constants}

\noindent DEMO\_USER\_ID: ID $=$ 00000000-0000-4000-8000-000000000001, the
proof of concept's one built-in user. It is fixed, so the seed script and
get\_current\_user always agree on who it is.

\medskip
\noindent DEFAULT\_TIME\_ZONE: String $=$ ``America/Toronto'', the team's own,
for a user made without one, such as the demo user. At Revision~0, sign-up
sends the phone's time zone instead.

\medskip
\noindent PLUGIN: Plugin $=$ (name: ``users'', router: None, lifespan: None,
health: None). Users has no routes yet. It is a plug-in anyway, so the sign-up
and sign-in routes of Revision~0 (FR1, FR4) are added here without changing
Control (M6).

\subsubsection{Exported Access Programs}

\begin{center}
\begin{tabular}{>{\raggedright\arraybackslash}p{3.6cm}
  >{\raggedright\arraybackslash}p{4.4cm}
  >{\raggedright\arraybackslash}p{2.4cm}
  >{\raggedright\arraybackslash}p{4.4cm}}
\toprule
\textbf{Name} & \textbf{In} & \textbf{Out} & \textbf{Exceptions} \\
\midrule
get\_current\_user & - & CurrentUser & NotAuthenticatedError \\
create\_user & Session, ID or None, String & ID & ValueError, IntegrityError \\
\bottomrule
\end{tabular}
\end{center}

\subsection{Semantics}

\subsubsection{State Variables}

\noindent users: set of tuple of (user\_id: ID, time\_zone: String,
created\_at: DateTime), every user, with each user\_id once. time\_zone is an
IANA time zone name. Nothing reads it yet. It is kept
because the server needs the user's date to refuse a daily total dated in the
future (SEC-24, HA~H51). The phone, not the server, puts each activity record
on its day (SEC-23).

\subsubsection{Environment Variables}

\noindent DEMO\_USER\_ENABLED: $\mathbb{B}$, read once per process, makes every
request that needs a user come from the demo user. It is off by default, so a
server not set up for the demo refuses them (SEC-15).

\subsubsection{Assumptions}

\begin{itemize}
\item The seed script (app/scripts/seed\_demo.py) makes the demo user, and its
  live portfolio in Wallet (M9), before a route needs them. get\_current\_user
  does not check that the user exists, so it needs no database query.
\end{itemize}

\subsubsection{Access Routine Semantics}

\noindent get\_current\_user():
\begin{itemize}
\item output: $out :=$ CurrentUser(DEMO\_USER\_ID). Routes ask for it by
  dependency injection, so the phone never sends a user ID, and a test can
  replace it to pick the user. At Revision~0 it finds the user from the session
  token instead, so no route changes.
\item exception: $exc :=$ ($\neg$DEMO\_USER\_ENABLED $\Rightarrow$
  NotAuthenticatedError)
\end{itemize}

{\raggedright
\noindent create\_user(session, user\_id, time\_zone). Unless they are given,
user\_id is None and time\_zone is DEFAULT\_TIME\_ZONE. A caller can give
user\_id so the seed script can make the demo user with DEMO\_USER\_ID.\par}
\begin{itemize}
\item transition: users $:=$ users $\cup \{(u$, time\_zone, the database's
  time$)\}$, where $u :=$ user\_id if it is given, else a new ID. It flushes,
  so the database makes the new ID, but never commits, so the caller can add
  the user's live portfolio and commit both together.
\item output: $out := u$
\item exception: $exc :=$ (time\_zone is not an IANA time zone name
  $\Rightarrow$ ValueError $|$ some user already has user\_id $\Rightarrow$
  IntegrityError). A bad time zone is refused before anything changes,
  since it would break finding the user's date (SEC-24).
\end{itemize}

\subsubsection{Local Functions}

None

\newpage

\section{MIS of Wallet} \label{MisWallet}

\subsection{Module}

Wallet (M9 in the MG)

\subsection{Uses}

Database (M1), Users (M8), Plugins (M22), JSON Format (M23)

\subsection{Syntax}

\subsubsection{Exported Types}

\noindent BalanceChangeReason $=$ \{activity\_reward, trade\_purchase\}, why a
balance changed. Every change keeps its reason (SEC-22), and the database
refuses any other. There is no reason for a sale yet, since the proof of
concept only buys. Selling (FR36) comes at Revision~0. Nothing pays a reward
yet either. That reason is there for Activity (M11), and tests use it to give
a portfolio currency to spend.

\medskip
{\raggedright
\noindent LockedPortfolio = tuple of (portfolio\_id: ID, user\_id: ID), the
read-only token that lock\_live\_portfolio returns (a lock token).
change\_balance asks for one, so changing a balance without locking it first is
hard to do by accident. Nothing checks where one came from.\par}

\medskip
{\raggedright
\noindent BalanceChangeReceipt = tuple of (balance\_change\_id: ID, amount:
Decimal, cash\_balance\_after: Decimal), a read-only copy of one change, so a
caller can't change the record through it (SEC-22). Trading (M12) keeps the change's ID with
each trade, so each trade points at the payment for it.\par}

\medskip
\noindent InsufficientFundsError, a DomainError with code
``insufficient\_funds'', status 409 and the message ``You don't have enough
currency for this.'' It refuses a purchase that costs more than the balance
(FR37). Status 409 says the request was understood but conflicts with the
balance now.

\subsubsection{Exported Constants}

\noindent PLUGIN: Plugin $=$ (name: ``wallet'', router: the route below,
lifespan: None, health: None)

\subsubsection{Exported Access Programs}

\begin{center}
% Small, with less space between columns, so names like
% NotAuthenticatedError fit their columns.
\small
\setlength{\tabcolsep}{3pt}
\begin{tabular}{>{\raggedright\arraybackslash}p{3.4cm}
  >{\raggedright\arraybackslash}p{3.9cm}
  >{\raggedright\arraybackslash}p{3.95cm}
  >{\raggedright\arraybackslash}p{4.0cm}}
\toprule
\textbf{Name} & \textbf{In} & \textbf{Out} & \textbf{Exceptions} \\
\midrule
lock\_live\_portfolio & Session, ID & LockedPortfolio & NoResultFound \\
change\_balance & Session, LockedPortfolio, Decimal, BalanceChangeReason, Date
or None & BalanceChangeReceipt & ValueError, InsufficientFundsError \\
create\_live\_portfolio & Session, ID & ID & IntegrityError \\
delete\_portfolio & Session, ID & - & - \\
live\_portfolio\_id & Session, ID & ID & NoResultFound \\
get\_balance & Session, ID & Decimal & NoResultFound \\
activity\_paid & Session, ID, Date & Decimal & ValueError \\
GET /v1/me/balance & - & tuple of (cash\_balance: Money) &
NotAuthenticatedError \\
\bottomrule
\end{tabular}
\end{center}

\subsection{Semantics}

\subsubsection{State Variables}

\noindent portfolios: set of tuple of (portfolio\_id: ID, user\_id: ID,
cash\_balance: Decimal, scenario\_id: ID or None, created\_at: DateTime). A
portfolio is live when its scenario\_id is None, and each user has at most one
live portfolio. Every cash\_balance is 0 or more (FR37).

\medskip
\noindent changes: set of tuple of (balance\_change\_id: ID, portfolio\_id: ID,
amount: Decimal, reason: BalanceChangeReason, activity\_date: Date or None,
created\_at: DateTime), every change to a balance, never edited. Each amount is
not 0, since a change of 0 would add a record without changing the balance. It
is above 0 exactly when the reason is activity\_reward, so a missing
minus sign can't turn a purchase into free currency, and the
change has an activity\_date, the day of activity it pays for, exactly when the
reason is activity\_reward.

\medskip
\noindent Every amount and balance has at most 4 decimals and 12 digits before
the point. Four decimals match Money (M23) and keep rounding far below
VALUE\_TOLERANCE (NFR1). For every portfolio $p$, $p$.cash\_balance is the sum of the amounts
of its changes (SEC-22). The database does not check this sum. change\_balance
keeps it true, since it is the only routine that changes a cash\_balance, and it
adds the change in the same transaction. Deleting a portfolio deletes its
changes, and other modules' rows that point at it with ON DELETE CASCADE, such
as the trades and holdings of Trading (M12).

\subsubsection{Environment Variables}

None

\subsubsection{Assumptions}

\begin{itemize}
\item A route that changes money calls lock\_live\_portfolio, then
  change\_balance, then commits once, by the commit rule of Database (M1). The
  lock lasts until the session's transaction ends, by a commit, a rollback or
  closing the session.
\item Every amount, and every balance a change leads to, is below $10^{12}$,
  the most the columns hold. A reward is at most the daily cap, so in practice
  none comes near it.
\item Every user a route acts for has a live portfolio. The seed script makes
  the demo user's.
\end{itemize}

\subsubsection{Access Routine Semantics}

\noindent lock\_live\_portfolio(session, user\_id), with $p$ the live portfolio
of user\_id:
\begin{itemize}
\item transition: $p$ is locked until session's transaction ends. Until
  then, a call for $p$ from another session waits, and it then sees what the
  first one saved. Taking turns keeps what a request reads after the lock
  true until it commits, such as whether its request ID was used
  (FR38) or what a day's rewards have paid (SEC-19, FR18).
\item output: $out :=$ LockedPortfolio($p$.portfolio\_id, user\_id)
\item exception: $exc :=$ (user\_id has no live portfolio $\Rightarrow$
  NoResultFound)
\end{itemize}

\noindent change\_balance(session, locked, amount, reason, activity\_date),
with activity\_date None by default and $p$ the portfolio with portfolio\_id
locked.portfolio\_id:
\begin{itemize}
\item transition: $p$.cash\_balance $:= p$.cash\_balance $+$ amount, and changes
  $:=$ changes $\cup \{(c$, $p$.portfolio\_id, amount, reason, activity\_date,
  the database's time$)\}$ with $c$ a new ID. Adding the change with the new
  balance keeps the balance equal to the sum of its changes (SEC-22). It
  flushes, so the database makes $c$, but never commits.
\item output: $out :=$ BalanceChangeReceipt($c$, amount, $p$.cash\_balance
  after the change)
\item exception: the first of these that holds is raised, before anything
  changes:\\
  \hspace*{1em}$exc :=$ (amount is 0, NaN or infinite, or has over 4 decimals
  $\Rightarrow$ ValueError\\
  \hspace*{1em}$|$ (reason $=$ activity\_reward) $\neq$ (amount $> 0$)
  $\Rightarrow$ ValueError\\
  \hspace*{1em}$|$ (reason $=$ activity\_reward) $\neq$ (activity\_date $\neq$
  None) $\Rightarrow$ ValueError\\
  \hspace*{1em}$|$ activity\_date is a DateTime $\Rightarrow$ ValueError\\
  \hspace*{1em}$|$ $p$.cash\_balance $+$ amount $< 0$ $\Rightarrow$
  InsufficientFundsError$)$\\
  A ValueError means the calling code is wrong, not a refusal the phone sees.
  An amount with over 4 decimals is refused, not rounded, since the database
  would round it quietly. A DateTime, which Python also counts as a Date, is
  refused, since its day depends on the time zone it is read in, so the same
  activity could be paid on two days (SEC-19). A balance never goes below 0, even if the caller skipped the lock,
  since the change and the check are one database statement.
\end{itemize}

\noindent create\_live\_portfolio(session, user\_id):
\begin{itemize}
\item transition: portfolios $:=$ portfolios $\cup \{(q$, user\_id, 0, None, the
  database's time$)\}$ with $q$ a new ID. It starts at 0, since it has no
  changes yet (SEC-22), and is kept apart from any scenario's (SEC-31). It flushes, so the database makes $q$, but never
  commits.
\item output: $out := q$
\item exception: $exc :=$ (user\_id is not in Users (M8) or already has a live
  portfolio $\Rightarrow$ IntegrityError)
\end{itemize}

\noindent delete\_portfolio(session, portfolio\_id):
\begin{itemize}
\item transition: the portfolio with portfolio\_id, its changes, and other
  modules' rows that point at it with ON DELETE CASCADE are removed in one
  statement, so nothing is left pointing at it. The seed script uses it to
  restart the demo. It never commits.
\item exception: None
\end{itemize}

\noindent live\_portfolio\_id(session, user\_id):
\begin{itemize}
\item output: $out :=$ the portfolio\_id of user\_id's live portfolio
\item exception: $exc :=$ (user\_id has no live portfolio $\Rightarrow$
  NoResultFound)
\end{itemize}

\noindent get\_balance(session, portfolio\_id):
\begin{itemize}
\item output: $out :=$ the cash\_balance of the portfolio with portfolio\_id,
  including changes session has made but not committed, so a route sees its
  own change before it commits
\item exception: $exc :=$ (no portfolio has portfolio\_id $\Rightarrow$
  NoResultFound)
\end{itemize}

\noindent activity\_paid(session, portfolio\_id, activity\_date):
\begin{itemize}
\item output: $out :=$ the sum of the amounts of the changes with
  portfolio\_id, reason \mbox{activity\_reward} and activity\_date, or 0.
  Nothing calls it yet. A step sync in Activity (M11) will use it to pay only
  the difference, so no activity is paid twice (SEC-19) or past the daily cap
  (FR18).
\item exception: $exc :=$ (activity\_date is a DateTime $\Rightarrow$
  ValueError), as in change\_balance
\end{itemize}

\noindent GET /v1/me/balance() takes session from get\_db of Database (M1)
and the user $u :=$ get\_current\_user() of Users (M8).
\begin{itemize}
\item output: $out :=$ (cash\_balance: $b$), where $b$ is get\_balance(session,
  $q$) as Money, and $q$ is live\_portfolio\_id(session, $u$.user\_id). The
  phone shows it, since rewards add to it (UC2) and a purchase can't cost more
  (FR37). It only reads, so it doesn't lock or commit.
\item exception: $exc :=$ the NotAuthenticatedError of get\_current\_user, with
  status 401
\end{itemize}

\subsubsection{Local Functions}

None

\newpage

\section{MIS of Market} \label{MisMarket}

\subsection{Module}

Market (M10 in the MG)

\subsection{Uses}

Database (M1), Clock (M2), Market Providers (M3), Constants (M7), Plugins (M22),
JSON Format (M23)

\subsection{Syntax}

\subsubsection{Exported Types}

\noindent AssetInfo = tuple of (asset\_id: ID, symbol: String, name: String,
is\_tradable: $\mathbb{B}$), a read-only copy of one asset, such as a stock.
is\_tradable is false once the asset is removed from the live market.

\medskip
\noindent CurrentPrice = tuple of (price: Decimal, quoted\_at: DateTime,
obtained\_at: DateTime, delayed: $\mathbb{B}$), an asset's newest price.
quoted\_at is when the provider set the price, and obtained\_at is when the
server got it.

\medskip
\noindent UnknownAssetError, a DomainError with code ``unknown\_asset'', status
404 and the message ``There is no asset with that symbol.''

\subsubsection{Exported Constants}

PLUGIN: Plugin $=$ (name: ``market'', router: the two routes below, lifespan:
lifespan, health: health). Both of Market's routes need no signed-in user. They
take now from get\_clock of Clock (M2) and session from get\_db of Database
(M1).

\subsubsection{Exported Access Programs}

\begin{center}
\begin{tabular}{>{\raggedright\arraybackslash}p{3.6cm}
  >{\raggedright\arraybackslash}p{2.4cm}
  >{\raggedright\arraybackslash}p{5.1cm} p{3.7cm}}
\toprule
\textbf{Name} & \textbf{In} & \textbf{Out} & \textbf{Exceptions} \\
\midrule
get\_asset\_by\_symbol & Session, String & AssetInfo & UnknownAssetError \\
get\_asset & Session, ID & AssetInfo & UnknownAssetError \\
get\_current\_price & Session, ID, DateTime & CurrentPrice or None & - \\
is\_delayed & DateTime, DateTime & $\mathbb{B}$ & - \\
is\_market\_open & DateTime & $\mathbb{B}$ & - \\
GET /v1/assets & - & tuple of (market\_open: $\mathbb{B}$, assets: seq of
tuple of (symbol: String, name: String, price: Money or None, price\_currency:
String, quoted\_at: DateTime or None, obtained\_at: DateTime or None, delayed:
$\mathbb{B}$)) & - \\
GET /v1/market/usage & - & tuple of (calls\_last\_minute: $\mathbb{Z}$,
busiest\_minute\_today: $\mathbb{Z}$, calls\_today: $\mathbb{Z}$,
limit\_per\_minute: $\mathbb{N}$) & - \\
\bottomrule
\end{tabular}
\end{center}

\subsection{Semantics}

\subsubsection{State Variables}

\noindent assets: set of tuple of (asset\_id: ID, symbol: String, name: String,
asset\_type: String, price\_currency: String, is\_tradable: $\mathbb{B}$), every
asset, with each symbol once.

\medskip
\noindent prices: set of tuple of (asset\_id: ID, price: Decimal,
price\_currency: String, quoted\_at: DateTime, obtained\_at: DateTime), every
price saved, at most one per asset and quoted\_at. An asset's current price is
the one with the latest quoted\_at. Every price\_currency, here and in assets, is
QUOTE\_CURRENCY from Constants (M7).

\medskip
\noindent requests: seq of tuple of (requested\_at: DateTime, symbol: String,
outcome: String), one per call to the provider in a committed poll (see
poll\_once).

\medskip
\noindent poller: \{off, running, stopped\}, the state of the price poller's
thread. It is off until lifespan starts the thread and again once the server
stops, running while the thread runs, and stopped if the thread ends while the
server runs.

\medskip
\noindent lastPollFailed: $\mathbb{B}$, whether the latest poll failed (see
run\_poller). It is false when lifespan starts.

\subsubsection{Environment Variables}

\noindent MARKET\_ALWAYS\_OPEN: $\mathbb{B}$, default false, treats the market as
always open.

\noindent MARKET\_POLLER\_ENABLED: $\mathbb{B}$, default false, starts the price
poller with the server.

\subsubsection{Assumptions}

\begin{itemize}
\item The database migrations add the assets. Market never adds or removes one.
\item At most one running server has MARKET\_POLLER\_ENABLED on.
\item Every time given to Market has a time zone.
\end{itemize}

\subsubsection{Access Routine Semantics}

\noindent get\_asset\_by\_symbol(session, symbol):
\begin{itemize}
\item output: $out := $ the AssetInfo of $a \in$ assets with $a$.symbol $=$
  symbol, even if not tradable
\item exception: $exc := $ (no $a \in$ assets has $a$.symbol $=$ symbol
  $\Rightarrow$ UnknownAssetError)
\end{itemize}

\noindent get\_asset(session, asset\_id):
\begin{itemize}
\item output: $out := $ the AssetInfo of $a \in$ assets with $a$.asset\_id $=$
  asset\_id
\item exception: $exc := $ (no $a \in$ assets has $a$.asset\_id $=$ asset\_id
  $\Rightarrow$ UnknownAssetError)
\end{itemize}

\noindent get\_current\_price(session, asset\_id, now):
\begin{itemize}
\item output: $out := $ None if the asset has no price yet, else (price,
  quoted\_at and obtained\_at of its current price $p$, delayed:
  is\_delayed($p$.quoted\_at, now))
\item exception: None
\end{itemize}

\noindent is\_delayed(quoted\_at, now):
\begin{itemize}
\item output: $out := $ is\_market\_open(now) $\land$ now $-$ quoted\_at $>$
  limit, where limit is QUOTE\_STALENESS\_LIMIT from Constants (M7)
\item exception: None
\end{itemize}

\noindent is\_market\_open(now):
\begin{itemize}
\item output: $out := $ MARKET\_ALWAYS\_OPEN $\lor$ is\_trading\_time(now)
\item exception: None
\end{itemize}

\noindent GET /v1/assets():
\begin{itemize}
\item output: $out := $ (market\_open: is\_market\_open(now), assets: for each
  tradable $a \in$ assets by symbol, its symbol, name and price\_currency, with
  the price (as Money), quoted\_at, obtained\_at and delayed from
  get\_current\_price, or None, None, None and false before its first price)
\item exception: None
\end{itemize}

\noindent GET /v1/market/usage():
\begin{itemize}
\item output: $out := $ (calls\_last\_minute: $|R(\text{now} - 1 \text{
  minute})|$, busiest\_minute\_today: the most entries of $R(\text{start})$ in
  one clock minute, or 0 if there are none, calls\_today:
  $|R(\text{start})|$, limit\_per\_minute: FINNHUB\_CALLS\_PER\_MINUTE of
  Market Providers (M3)), where $R(t)$ is the entries of requests from $t$ to
  now, and start is the midnight that begins now's New York day
\item exception: None
\end{itemize}

\subsubsection{Local Functions}

\noindent HOLIDAYS: set of Date, the weekdays in 2026 and 2027 that the New
York Stock Exchange closes for a holiday.\\
EARLY\_CLOSES: set of Date = \{2026-11-27, 2026-12-24, 2027-11-26\}, its early
closes.

\smallskip
\noindent is\_trading\_time: DateTime $\rightarrow \mathbb{B}$\\
is\_trading\_time(now) $:= d$ is a weekday not in HOLIDAYS $\land\ 9{:}30 \leq
t < c$, where $d$ and $t$ are now's day and time in New York (NEW\_YORK, the
America/New\_York time zone), and $c$ is 13:00 if $d \in$ EARLY\_CLOSES, else
16:00.

\smallskip
\noindent POLL\_INTERVAL: Duration = 60 seconds, the wait after each poll.\\
POLL\_AFTER\_CLOSE: Duration = 15 minutes, how long polling goes on after the
close.\\
STOP\_TIMEOUT: Duration = 15 seconds, how long the server waits for the poller
to stop.

\smallskip
\noindent should\_poll: DateTime $\rightarrow \mathbb{B}$\\
should\_poll(now) $:=$ is\_market\_open(now) $\lor$ is\_market\_open(now $-$
POLL\_AFTER\_CLOSE)

\smallskip
\noindent poll\_once(session: Session, provider: MarketDataProvider, now:
DateTime): for each tradable asset, it gets a quote from provider.get\_quote,
saves it with save\_quote and adds (now, symbol, outcome) to requests. outcome
is ``ok'' for a quote. A RateLimitedError gives ``rate\_limited'' and ends the
poll, an UnknownSymbolError gives ``unknown\_symbol'' and any other
MarketDataError gives ``error''. It does not commit. It returns true only if
it saved a quote for every tradable asset.

\smallskip
\noindent save\_quote(session: Session, asset\_id: ID, $q$: Quote) adds $q$ to
prices with obtained\_at the database's time. If the asset has a price at
$q$.quoted\_at already, an equal price is kept and a different one is replaced,
with a new obtained\_at.

\smallskip
\noindent run\_poller(stop: $\mathbb{B}$, provider: MarketDataProvider): until
stop is true, if should\_poll(now) with now $:=$ system\_clock(), it runs
poll\_once at now in a new SessionLocal session and commits, then waits
POLL\_INTERVAL. A poll that raises an error is logged and rolled back. After
each poll, lastPollFailed $:=$ (the poll raised an error, or poll\_once
returned false). So a rate limit, a provider that is down, a wrong key or a
single symbol without a price all count as a failed poll.

\smallskip
\noindent lifespan: Lifespan. At start it reads the two environment variables
and get\_provider\_settings of Market Providers (M3) and runs
check\_always\_open, and a ValidationError or ValueError stops the server. If
MARKET\_POLLER\_ENABLED, run\_poller runs on its own thread with the provider
from build\_provider until the server stops, which waits up to STOP\_TIMEOUT.
It sets poller and lastPollFailed as their descriptions say.

\smallskip
\noindent health: HealthCheck\\
health() $:=$ (poller $=$ off $\Rightarrow$ (``ok'', ``poller off'') $|$
poller $=$ stopped $\Rightarrow$ (\mbox{``degraded''}, ``poller stopped'') $|$
poller $=$ running $\land$ lastPollFailed $\Rightarrow$ (\mbox{``degraded''},
``last poll failed'') $|$ poller $=$ running $\land \neg$lastPollFailed
$\Rightarrow$ (``ok'', None)). A failed poll shows until the next poll, which
may be the next trading day.

\smallskip
\noindent check\_always\_open(market, providers: ProviderSettings) raises
ValueError if market has MARKET\_ALWAYS\_OPEN on and providers.market\_provider
$\neq$ ``fake''.

\newpage

% M11 Activity: #249
\section{MIS of Trading} \label{MisTrading}

\subsection{Module}

Trading (M12 in the MG)

\subsection{Uses}

Database (M1), Clock (M2), Users (M8), Wallet (M9), Market (M10), Plugins
(M22), JSON Format (M23)

\subsection{Syntax}

\subsubsection{Exported Types}

{\raggedright\noindent Trading's errors are DomainErrors with HTTP status
409: RequestIdReusedError (code request\_id\_reused), MarketClosedError
(market\_closed), PriceOutOfDateError (price\_out\_of\_date) and
PriceChangedError (price\_changed). The local function buy says when each is
raised.\par}

\medskip
\noindent HoldingValue = tuple of (symbol: String, quantity: $\mathbb{Z}$,
price: Money, quoted\_at: DateTime, delayed: $\mathbb{B}$, value: Money), one
holding in the reply to GET /v1/me/portfolio.

\subsubsection{Exported Constants}

PLUGIN: Plugin $=$ (name: ``trading'', router: the two routes below, lifespan:
None, health: None)

\subsubsection{Exported Access Programs}

\begin{center}
\begin{tabular}{>{\raggedright}p{3.2cm} >{\raggedright}p{2.1cm}
  >{\raggedright}p{5.2cm} >{\raggedright\arraybackslash}p{4.3cm}}
\toprule
\textbf{Name} & \textbf{In} & \textbf{Out} & \textbf{Exceptions} \\
\midrule
POST /v1/me/trades & ID, String, String, $\mathbb{N}$, Money
  & tuple of (trade\_id: ID, status: String, unit\_price: Money,
  currency\_amount: Money, executed\_at: DateTime, cash\_balance: Money)
  & RequestIdReusedError, MarketClosedError, UnknownAssetError,
  PriceOutOfDateError, PriceChangedError, InsufficientFundsError \\
GET /v1/me/portfolio & - & tuple of (cash\_balance: Money, holdings: seq of
  HoldingValue, holdings\_value: Money, total\_value: Money, simulated:
  $\mathbb{B}$) & - \\
\bottomrule
\end{tabular}
\end{center}

\subsection{Semantics}

\subsubsection{State Variables}

\noindent trades: set of Trade, every trade made, where Trade = tuple of
(trade\_id: ID, request\_id: ID, portfolio\_id: ID, asset\_id: ID,
trade\_type: String, quantity: $\mathbb{N}$, unit\_price: Decimal,
currency\_amount: Decimal, cash\_balance\_after: Decimal, status: String,
balance\_change\_id: ID, created\_at: DateTime). A request\_id appears at most
once among a portfolio's trades. Every trade is a ``buy'' with status
``executed''.

\medskip
\noindent holdings: set of tuple of (portfolio\_id: ID, asset\_id: ID,
quantity: $\mathbb{Z}$), the shares each portfolio holds, at most one entry per
portfolio and stock, with quantity never below 0.

\medskip
\noindent When Wallet deletes a portfolio, its trades and holdings go with it.

\subsubsection{Environment Variables}

None

\subsubsection{Assumptions}

The lock from lock\_live\_portfolio lasts until the session commits or is
closed, so requests that change one portfolio run one at a time. Every stock in
holdings has at least one price.

\subsubsection{Access Routine Semantics}

\noindent Both routes take session from get\_db, now from get\_clock and $u :=$
get\_current\_user(), which refuses with status 401 when there is no current
user.

\medskip
\noindent POST /v1/me/trades(request\_id, symbol, trade\_type, quantity,
expected\_price):
\begin{itemize}
\item transition: the route locks the portfolio with locked $:=$
  lock\_live\_portfolio(session, $u$.user\_id), runs $(t, \mathit{new}) :=$
  buy(session, locked, request, now) on the request body, and commits once. So
  the payment, the shares and the trade are saved together or not at all.
  With the lock first, two copies of one request run one after the other, and
  if the first made a trade, the second finds it.
\item output: the status is 201 for a new trade and 200 for a repeat, and
  $out :=$ ($t$.trade\_id, $t$.status, $t$.unit\_price,
  $t$.currency\_amount, $t$.created\_at, $t$.cash\_balance\_after), so a
  repeated request\_id gets the same body as the first answer.
\item exception: $exc :=$ the exception buy raises, with status 404 for
  UnknownAssetError and 409 for the rest. A refused request saves nothing, so
  sending it again runs every check again. A trade\_type other than ``buy'' gets status 422.
\end{itemize}

\noindent GET /v1/me/portfolio():
\begin{itemize}
\item output: $out :=$ portfolio\_value(session, live\_portfolio\_id(session,
  $u$.user\_id), now) with simulated $=$ true, so the phone labels the values
  as simulated.
\item exception: None
\end{itemize}

\subsubsection{Local Functions}

\noindent TradeRequest = tuple of (request\_id: ID, symbol: String,
trade\_type: String, quantity: $\mathbb{N}$, expected\_price: Money), the body
of POST /v1/me/trades.

\medskip
\noindent buy(session: Session, locked: LockedPortfolio, request: TradeRequest,
now: DateTime) $\rightarrow$ tuple of (t: Trade, new: $\mathbb{B}$) holds the
trade rules and never commits. Here $p :=$ locked.portfolio\_id, and old is the
trade in trades with portfolio\_id $p$ and request\_id request.request\_id, or
None. old matches request when it has the same symbol and quantity, and
old.unit\_price $=$ request.expected\_price. For a new request (old $=$ None),
$a :=$ get\_asset\_by\_symbol(session, request.symbol), $c :=$
get\_current\_price(session, $a$.asset\_id, now) and cost $:=$
request.quantity $\times$ $c$.price.
\begin{itemize}
\item transition: for a new request that passes every check, $b :=$
  change\_balance(session, locked, $-$cost, trade\_purchase) pays for it, $p$'s
  holding of $a$ gains request.quantity shares (from 0 if none), and trades
  $:=$ trades $\cup \{t\}$, where $t :=$ (a new ID, request.request\_id, $p$,
  $a$.asset\_id, ``buy'', request.quantity, $c$.price, cost,
  get\_balance(session, $p$) after the payment, ``executed'',
  $b$.balance\_change\_id, now).
\item output: $out :=$ (old matches request $\Rightarrow$ (old, false) $\mid$
  old $=$ None $\Rightarrow$ ($t$, true)), so a repeat gets its first trade
  back, even after the market closes or the price moves.
\item exception: for a repeat, $exc := ($old $\neq$ None $\land$ old does not
  match request $\Rightarrow$\\
  \hspace*{1em}RequestIdReusedError$)$.\\
  For a new request the first of these
  that holds is raised, in this order:\\
  \hspace*{1em}$exc := (\neg$is\_market\_open(now) $\Rightarrow$
    MarketClosedError\\
  \hspace*{1em}$\mid$ no asset has symbol request.symbol $\Rightarrow$
    UnknownAssetError\\
  \hspace*{1em}$\mid c =$ None $\lor$ $c$.delayed $\Rightarrow$
    PriceOutOfDateError\\
  \hspace*{1em}$\mid c$.price $\neq$ request.expected\_price $\Rightarrow$
    PriceChangedError\\
  \hspace*{1em}$\mid$ get\_balance(session, $p$) $<$ cost $\Rightarrow$
    InsufficientFundsError$)$\\
  PriceChangedError's message gives the new price, $c$.price.
\end{itemize}

\noindent portfolio\_value(session: Session, portfolio\_id: ID, now: DateTime)
$\rightarrow$ tuple of (cash\_balance: Decimal, holdings: seq of HoldingValue,
holdings\_value: Decimal, total\_value: Decimal) only reads.
\begin{itemize}
\item output: $out := (b, H, s, b + s)$, where $b :=$ get\_balance(session,
  portfolio\_id), $H$ has one entry per holding $h$ of portfolio\_id, in order
  of symbol, and $s$ is the sum of their values. With $c :=$
  get\_current\_price(session, $h$.asset\_id, now), $h$'s entry is (its symbol,
  $h$.quantity, $c$.price, $c$.quoted\_at, $c$.delayed, $h$.quantity $\times$
  $c$.price). A delayed price is still used, and delayed flags it.
\item exception: None
\end{itemize}

\newpage

% M13 to M17 and M21 are built after the PoC and are specified at Revision 0.
% M18 Navigation, M19 Screens and M20 Sync Service (app): #250
\section{MIS of Plugins} \label{MisPlugins}

\subsection{Module}

Plugins (M22 in the MG)

\subsection{Uses}

None

\subsection{Syntax}

\subsubsection{Exported Types}

\noindent Plugin = tuple of (name: String, router: Router or None, lifespan:
Lifespan or None, health: HealthCheck or None), a record of what one module adds
to the server. It cannot be changed once built.
\begin{itemize}
\item name: the module's folder name, such as ``market''.
\item router: the module's routes, which Control (M6) adds under /v1.
\item lifespan: the module's background work, which runs while the server runs.
\item health: the module's health check, which Control (M6) shows in
  GET~/v1/health.
\end{itemize}

\noindent Lifespan = App $\rightarrow$ work that starts with the server and
stops when it shuts down, such as the price poller of Market (M10). A Lifespan
stops its work in a finally block, so the work still stops if another module
fails to start.

\medskip
\noindent HealthCheck = () $\rightarrow$ ModuleHealth, a check of the
module's own state.

\medskip
\noindent ModuleHealth = tuple of (status: String, detail: String or None), a
module's health. Its status is ``ok'' or \mbox{``degraded''}, and its detail is a
short fixed text for the team, such as ``poller stopped''.

\medskip
\noindent Router: a group of HTTP routes (a FastAPI APIRouter).

\medskip
\noindent App: the web app that serves the routes (a FastAPI app).

\subsubsection{Exported Constants}

None

\subsubsection{Exported Access Programs}

\begin{center}
\begin{tabular}{p{2.4cm} p{6.4cm} p{2cm} p{2.6cm}}
\toprule
\textbf{Name} & \textbf{In} & \textbf{Out} & \textbf{Exceptions} \\
\midrule
new Plugin & String, Router or None, Lifespan or None, HealthCheck or None &
  Plugin & ValueError \\
\bottomrule
\end{tabular}
\end{center}

\subsection{Semantics}

\subsubsection{State Variables}

None

\subsubsection{Environment Variables}

None

\subsubsection{Assumptions}

\begin{itemize}
\item Every Lifespan follows the start and stop rule in its type above.
\item A HealthCheck reads only its module's own state in memory, with no
  database or network call, so it answers at once. Its detail never holds a
  secret or an error's own text, because GET /v1/health needs no signed-in
  user.
\end{itemize}

\subsubsection{Access Routine Semantics}

\noindent new Plugin(name, router, lifespan, health):
\begin{itemize}
\item output: $out :=$ (name: name, router: router, lifespan: lifespan, health:
  health). Each of router, lifespan and health is None when the module leaves
  it out.
\item exception: $exc :=$ (lifespan $\neq$ None $\land$ health $=$ None
  $\Rightarrow$ ValueError). A module with background work must report its
  health, because that work can fail without any request seeing it.
\end{itemize}

\subsubsection{Local Functions}

None

\newpage

\section{MIS of JSON Format} \label{MisJSON}

\subsection{Module}

JSON Format (M23 in the MG)

\subsection{Uses}

None

\subsection{Syntax}

\subsubsection{Exported Types}

\noindent Money = Decimal, an amount, such as a price. It is sent to the phone
as a JSON string with exactly four decimals (see \_to\_text).

\medskip
\noindent DomainError = tuple of (code: String, message: String, status:
$\mathbb{N}$), an error a route raises to refuse a request it understood. The
code names the kind of refusal, the message says why in words a user can read,
and the status is the HTTP status of the reply. The defaults are code
``domain\_error'', message ``This request could not be completed.'' and status
400. Each kind of refusal is a subtype of DomainError with its own code, and it
can also set its own message and status.

\medskip
\noindent Refusal = tuple of (detail: tuple of (code: String, message:
String)), the JSON body Control (M6) sends, with HTTP status $e$.status, for
every DomainError $e$ a route raises:

\begin{center}
\verb|{"detail": {"code": e.code, "message": e.message}}|
\end{center}

\noindent FastAPI itself refuses a request that does not fit a route's types
with status 422, in its own shape.

\subsubsection{Exported Constants}

None

\subsubsection{Exported Access Programs}

\begin{center}
\begin{tabular}{p{3.4cm} p{4.6cm} p{2.4cm} p{2.9cm}}
\toprule
\textbf{Name} & \textbf{In} & \textbf{Out} & \textbf{Exceptions} \\
\midrule
new Money & Decimal, $\mathbb{Z}$, $\mathbb{R}$ or String & Money &
  ValidationError \\
new DomainError & String or None & DomainError & - \\
\bottomrule
\end{tabular}
\end{center}

\subsection{Semantics}

\subsubsection{State Variables}

None

\subsubsection{Environment Variables}

None

\subsubsection{Assumptions}

Every amount a route sends to the phone has the type Money. An amount that
comes from a division or an average is rounded to four decimals before it goes
in a Money field.

\subsubsection{Access Routine Semantics}

\noindent new Money(amount):
\begin{itemize}
\item output: $out := $ the Decimal that amount names
\item exception: $exc := $ (amount is not a finite number, or cannot be
  written exactly with four decimals $\Rightarrow$ ValidationError). It is
  refused, never rounded.
\end{itemize}

\noindent new DomainError(message):
\begin{itemize}
\item output: $out :=$ (code: $k$.code, message: (message $\neq$ None
  $\Rightarrow$ message $|$ message $=$ None $\Rightarrow$ $k$.message),
  status: $k$.status), where $k$ is the kind of DomainError being made and
  message is None when left out
\item exception: None
\end{itemize}

\subsubsection{Local Functions}

\noindent \_to\_text: Money $\rightarrow$ String

\noindent \_to\_text($m$) $:=$ $m$ with exactly four digits after the decimal
point, such as ``84.0000'' for 84. A zero is always ``0.0000'', never
``-0.0000''.

\newpage

% M24 API Service (app): #250

\bibliographystyle {plainnat}
\bibliography {../../../refs/References}

\newpage

\newpage{}

\section{Appendix --- Generative AI Use Disclosure}

\wss{ChatGPT (version 5) was used to brainstorm ideas for the template for this disclose section.}

\subsection{AI Tools Used}

\wss{Give the name of the AI tool(s) used and the version.}

\subsection{How and Where Used}

\wss{Explain what the tool(s) were used for.  Examples include: brainstorming,
explaining a technical concept, reviewing the document, generating or modifying
text, generating or modifying diagrams, identifying errors or omissions or
inconsistencies.}

\wss{For each of the uses, identify which parts of the document were affected.}

\wss{You do not need to report every interaction, but you should disclose the 
uses that materially contributed to your work}

\subsection{Verification of AI-Generated Content}

\wss{\begin{itemize}
    \item Describe how the team checked AI-generated or AI-assisted material
    before including it in the document.
    \item In particular, verify factual claims, technical assertions,
    requirements, calculations, references, and other information for which an
    AI system may produce incorrect or fabricated results.
    \item \textbf{The team is responsible for the accuracy, completeness,
consistency, and appropriateness of the submitted document, regardless of
whether material was generated by a human or an AI system.}
\end{itemize}}

\newpage{}

\section{Appendix --- Reflection}

\wss{Not required for CAS 741 projects}

The information in this section will be used to evaluate the team members on the
graduate attribute of Problem Analysis and Design.

\input{../../Reflection.text}

\begin{enumerate}
  \item What went well while writing this deliverable? 
  \item What pain points did you experience during this deliverable, and how
    did you resolve them?
  \item Which of your design decisions stemmed from speaking to your client(s)
  or a proxy (e.g. your peers, stakeholders, potential users)? For those that
  were not, why, and where did they come from?
  \item While creating the design doc, what parts of your other documents (e.g.
  requirements, hazard analysis, etc), it any, needed to be changed, and why?
  \item What are the limitations of your solution?  Put another way, given
  unlimited resources, what could you do to make the project better? (LO\_ProbSolutions)
  \item Give a brief overview of other design solutions you considered.  What
  are the benefits and tradeoffs of those other designs compared with the chosen
  design?  From all the potential options, why did you select the documented design?
  (LO\_Explores)
  \item (After you have implemented another team's module, which means this
  isn't filled in until after the original deadline). What did you learn by
  implementing another team's module? Were all the details you needed in the
  documentation, or did you need to make assumptions, or ask the other team
  questions?  If your team also had another team implement one of your modules,
  what was this experience like?  Are there things in your documentation you
  could have changed to make the process go more smoothly for when an
  ``outsider'' completes some of the implementation?

  Your reflection should be about concepts learned in your SE education, such as modular decomposition, cohesion, coupling, information hiding, etc.  You should critique the quality and utility of the specification you received.

  You should create a PR for your implemented module.  A link to the PR should be part of your reflection.

\end{enumerate}

\end{document}